% European Heart Journal Issue @ A Glance ACS 2026 教學講義
% 急性冠心症專題：機轉、診斷、心因性休克處置與新治療標靶
\documentclass[12pt,a4paper]{article}
% Drake_preamble.tex - LaTeX Preamble Template
% 作者: 謝慕揚, MD, PhD, FESC
% 
% 使用說明:
% 1. 表格請使用 longtable，不要有垂直分隔線 |
% 2. 表格內換行使用 \newline，不要使用 \\
% 3. 藥物名稱、檢驗、檢查請維持英文
% 4. Reference 格式請使用 \href 加上驗證過的 DOI

% Including essential packages for Chinese typesetting
\usepackage{xeCJK}
\usepackage{fontspec}
% Configuring Chinese font with Noto Serif CJK TC, fallback to SimSun
\IfFontExistsTF{Noto Serif CJK TC}{
  \setCJKmainfont{Noto Serif CJK TC}
  \setCJKsansfont{Noto Serif CJK TC}
  \setCJKmonofont{Noto Serif CJK TC}
}{\setCJKmainfont{SimSun}
  \setCJKsansfont{SimSun}
  \setCJKmonofont{SimSun}
}

% Including other necessary packages
\usepackage{geometry}
\usepackage{titlesec}
\usepackage{enumitem}
\usepackage{xcolor}
\usepackage{colortbl}
\usepackage{tcolorbox}
\usepackage{booktabs}
\usepackage{longtable}
\usepackage{array}
\usepackage{graphicx}
\usepackage{tikz}
\usepackage{fancyhdr}
\usepackage{multicol}
\usepackage{datetime2}
\usepackage{hyperref}
\usepackage{mdframed}
\usepackage{amsmath}
\usepackage{amssymb}
\usepackage{multirow}

\hypersetup{
    colorlinks=true,
    linkcolor=emergency_blue,
    citecolor=emergency_blue,
    urlcolor=emergency_blue,
    pdfborder={0 0 0}
}

% Defining page geometry
\geometry{
  top=2.5cm,
  bottom=2.5cm,
  left=2cm,
  right=2cm,
  headheight=25pt
}

% Adjusting topmargin to compensate for increased headheight
\addtolength{\topmargin}{-10pt}

% Defining colors
\definecolor{emergency_red}{RGB}{186,24,27}
\definecolor{emergency_blue}{RGB}{0,114,188}
\definecolor{emergency_gray}{RGB}{120,120,120}
\definecolor{highlight_yellow}{RGB}{255,250,205}

% Setting title formats
\titleformat{\section}[block]{\Large\bfseries\color{emergency_red}}{\thesection}{10pt}{}
\titleformat{\subsection}[block]{\large\bfseries\color{emergency_blue}}{\thesubsection}{8pt}{}
\titleformat{\subsubsection}[block]{\normalsize\bfseries\color{emergency_gray}}{\thesubsubsection}{6pt}{}

% Defining custom environment for important notes
\newmdenv[
    linecolor=emergency_red,
    backgroundcolor=highlight_yellow,
    linewidth=2pt,
    topline=true,
    bottomline=true,
    leftline=true,
    rightline=true,
    innertopmargin=10pt,
    innerbottommargin=10pt,
    innerleftmargin=10pt,
    innerrightmargin=10pt
]{importantbox}

% Defining note command with tcolorbox
\newcommand{\note}[1]{
    \par
    \noindent
    \begin{tcolorbox}[
        colback=emergency_blue!5!white,
        colframe=emergency_blue,
        title=\textbf{重點提醒},
        fonttitle=\bfseries,
        boxrule=0.5pt,
        arc=3pt,
        left=6pt,
        right=6pt,
        top=6pt,
        bottom=6pt
    ]
    #1
    \end{tcolorbox}
}

% Key findings box
\newcommand{\keyfinding}[1]{
    \par
    \noindent
    \begin{tcolorbox}[
        colback=yellow!10!white,
        colframe=orange!75!black,
        title=\textbf{核心發現摘要},
        fonttitle=\bfseries,
        boxrule=0.5pt,
        arc=3pt
    ]
    #1
    \end{tcolorbox}
}

% Defining custom date format for ISO 8601
\DTMnewdatestyle{iso}{
  \renewcommand{\DTMdisplaydate}[4]{\number##1-\number##2-\number##3}
  \renewcommand{\DTMDisplaydate}[4]{\number##1-\number##2-\number##3}
}
\DTMsetdatestyle{iso}
\newcommand{\mydate}{\DTMtoday}


% Custom header for this document
\pagestyle{fancy}
\fancyhf{}
\fancyhead[L]{\textcolor{emergency_red}{European Heart Journal 教學講義}}
\fancyhead[R]{\textcolor{emergency_blue}{ACS Issue @ A Glance}}
\fancyfoot[C]{\textcolor{emergency_gray}{\thepage}}
\renewcommand{\headrulewidth}{1pt}
\renewcommand{\headrule}{\hbox to\headwidth{\color{emergency_red}\leaders\hrule height \headrulewidth\hfill}}

\begin{document}

%============================================================
% Title Page
%============================================================
\begin{center}
{\LARGE\bfseries\textcolor{emergency_red}{European Heart Journal}}\\[0.5cm]
{\Large\textbf{Issue @ A Glance: Focus on Acute Coronary Syndrome}}\\[0.3cm]
{\large 期刊焦點：急性冠心症——機轉、診斷、心因性休克處置與新治療標靶}\\[0.5cm]
{\normalsize \href{https://doi.org/10.1093/eurheartj/ehag016}{Eur Heart J 2026;47:387-392. DOI: 10.1093/eurheartj/ehag016}}\\[0.3cm]
{\normalsize 整理：謝慕揚 MD, PhD, FESC \quad 日期：\mydate}
\end{center}

\vspace{0.5cm}

\keyfinding{
本期 EHJ 聚焦急性冠心症（ACS），涵蓋六大主題：
\begin{enumerate}[noitemsep]
    \item PCI 後心絞痛的病因與處置
    \item Non-STEMI ACS 的完整光譜與管理
    \item 心肌梗塞通用定義的爭議
    \item Clonal haematopoiesis（CHIP）與冠心病死亡率
    \item 心因性休克的 Impella 使用（J-PVAD Registry）
    \item 長型 Troponin T 在 MI 診斷的應用
\end{enumerate}
}

%============================================================
\section{PCI 後心絞痛}
%============================================================

\subsection{問題概述}

\begin{itemize}
    \item PCI 是全球最主要的心肌血管重建方式
    \item PCI 後心絞痛持續或復發是常見問題
    \item 違背穩定冠心病患者接受 PCI 的主要目的：緩解症狀
    \item 現行指引對此問題的診斷、預防與處置著墨有限
\end{itemize}

\subsection{可能原因}

\begin{longtable}{p{4cm}p{10cm}}
\toprule
\textbf{類別} & \textbf{原因} \\
\midrule
\endfirsthead
\midrule
\textbf{類別} & \textbf{原因} \\
\midrule
\endhead
\bottomrule
\endfoot
PCI 結果次優 & • 殘餘狹窄\newline • 支架 underexpansion\newline • 支架 malapposition\newline • 邊緣夾層\newline • 功能性結果不佳 \\
\midrule
非阻塞性病因 & • 冠狀動脈微血管疾病（CMD）\newline • 冠狀動脈痙攣\newline • 心肌橋 \\
\midrule
進展性疾病 & • 新發病灶\newline • 支架內再狹窄\newline • 支架血栓 \\
\bottomrule
\end{longtable}

\subsection{建議處置流程}

\begin{enumerate}
    \item \textbf{詳細評估}：區分心因性 vs 非心因性胸痛
    \item \textbf{功能性冠狀動脈造影}：利用新工具評估
    \begin{itemize}
        \item FFR/iFR 評估狹窄血流動力學意義
        \item CFR/IMR 評估微血管功能
        \item Acetylcholine test 評估血管痙攣
    \end{itemize}
    \item \textbf{影像評估}：IVUS/OCT 確認支架相關問題
    \item \textbf{個人化治療}：針對病因給予適當處置
\end{enumerate}

%============================================================
\section{Non-STEMI ACS 的完整光譜}
%============================================================

\subsection{定義與分類}

\begin{itemize}
    \item \textbf{ACS}：阻塞性冠狀動脈血栓併發缺血性心肌疾病
    \item 依 ECG 表現分類：
    \begin{itemize}
        \item STEMI：ST 段上升心肌梗塞
        \item NSTE-ACS：非 ST 段上升 ACS（含或不含 troponin 上升）
    \end{itemize}
\end{itemize}

\subsection{初步診斷與風險分層}

\begin{longtable}{p{3cm}p{11cm}}
\toprule
\textbf{要素} & \textbf{內容} \\
\midrule
\endfirsthead
\bottomrule
\endfoot
症狀 & 典型胸痛特徵、持續時間、誘發因素 \\
\midrule
ECG 變化 & ST 下降、T 波倒置、動態變化 \\
\midrule
Troponin & 上升且有動態變化（rise and/or fall pattern）\\
\midrule
排除診斷 & • Type 2 MI（供需失衡）\newline • 非缺血性心肌損傷（myocarditis、takotsubo 等）\newline • 全身性疾病導致的心肌損傷 \\
\bottomrule
\end{longtable}

\subsection{冠狀動脈造影的角色}

\begin{itemize}
    \item 評估是否有阻塞性冠心病
    \item 區分 MINOCA（MI with Non-Obstructive Coronary Arteries）
    \item 評估是否適合血管重建
    \item \textbf{完整血管重建}：多血管疾病患者，完整血管重建較僅處理罪犯血管可降低死亡率
\end{itemize}

\subsection{後續治療原則}

\begin{itemize}
    \item \textbf{個人化抗血栓治療}：依出血與缺血風險調整
    \item \textbf{疾病調修二級預防}：
    \begin{itemize}
        \item 高強度 statin
        \item RAAS 抑制劑
        \item Beta-blocker
        \item 生活型態介入
    \end{itemize}
\end{itemize}

%============================================================
\section{心肌梗塞通用定義的爭議}
%============================================================

\subsection{第四版通用定義（4th UDMI）重點}

\begin{itemize}
    \item \textbf{病理定義}：長時間缺血導致的心肌細胞死亡
    \item \textbf{臨床定義}：缺血臨床證據 + 心肌細胞死亡生化標記
    \item \textbf{Cardiac troponin}：標準生化標記
    \item \textbf{MI 分類}：依病理生理機轉分為 Type 1-5
\end{itemize}

\subsection{Type 1 vs Type 2 MI}

\begin{longtable}{p{2cm}p{6cm}p{6cm}}
\toprule
& \textbf{Type 1 MI} & \textbf{Type 2 MI} \\
\midrule
\endfirsthead
\bottomrule
\endfoot
機轉 & 動脈粥狀硬化斑塊破裂/侵蝕導致血栓 & 心肌氧氣供需失衡 \\
\midrule
病因 & 自發性冠狀動脈事件 & 心律不整、貧血、低血壓、低血氧等 \\
\midrule
造影 & 通常有罪犯血管 & 可能無明顯阻塞或多處狹窄 \\
\midrule
處置 & 冠狀動脈介入 & 處理根本原因 \\
\bottomrule
\end{longtable}

\subsection{爭議焦點}

\begin{itemize}
    \item 通用定義旨在\textbf{定義 MI}，而非提供治療指引
    \item 臨床實務中區分 Type 1 vs Type 2 MI 仍有困難
    \item 非缺血性心肌損傷（myocardial injury）的分類
    \item 高敏感度 troponin 時代的挑戰
\end{itemize}

%============================================================
\section{CHIP 與冠心病死亡率}
%============================================================

\subsection{什麼是 CHIP？}

\begin{itemize}
    \item \textbf{Clonal Haematopoiesis of Indeterminate Potential}
    \item 造血幹細胞獲得體細胞突變，形成優勢克隆
    \item 隨年齡增加而常見
    \item 與心血管風險相關，但機轉尚未完全釐清
\end{itemize}

\subsection{研究發現}

\begin{longtable}{p{4cm}p{10cm}}
\toprule
\textbf{項目} & \textbf{內容} \\
\midrule
\endfirsthead
\midrule
\textbf{項目} & \textbf{內容} \\
\midrule
\endhead
\bottomrule
\endfoot
研究方法 & 8,612 位冠狀動脈造影確診 CAD 患者\newline 13 個 CHIP 驅動基因定序 \\
\midrule
配對分析 & CHIP 帶原者與非帶原者 1:1 傾向分數配對（2,389 對）\\
\midrule
主要發現 & CHIP 增加 3 年死亡率（HR 1.39, P < .001）\\
\midrule
相關基因 & TET2、ASXL1、DNMT3A、JAK2、PPM1D、SF3B1、SRSF2、U2AF1 \\
\bottomrule
\end{longtable}

\subsection{TET2 突變的機轉研究}

\begin{itemize}
    \item TET2 CHIP 帶原者動脈斑塊特徵：
    \begin{itemize}
        \item 壞死核心增大
        \item 發炎增加
        \item 斑塊穩定性降低
    \end{itemize}
    \item Multi-omics 分析：脂質代謝與發炎路徑上調
    \item TET2\textsuperscript{+/-} 巨噬細胞：
    \begin{itemize}
        \item LDLR 表現增加
        \item 脂質攝取增加
        \item 與 LDLR 啟動子染色質可及性增加相關
    \end{itemize}
\end{itemize}

\begin{importantbox}
\textbf{臨床意義：}\\
CHIP 是 CAD 患者死亡率的獨立預測因子。TET2 突變透過 LDLR 上調與發炎活化促進促動脈粥狀硬化巨噬細胞表型，連結表觀遺傳失調與 CAD 不良預後。
\end{importantbox}

%============================================================
\section{心因性休克的 Impella 使用：J-PVAD Registry}
%============================================================

\subsection{背景：DanGer Shock Trial}

\begin{itemize}
    \item DanGer Shock 試驗顯示 Impella 在\textbf{高度篩選}的 STEMI 相關心因性休克患者有生存獲益
    \item 但真實世界中，Impella 常用於不符合試驗標準的患者
\end{itemize}

\subsection{DanGer Shock 納入標準}

\begin{enumerate}
    \item STEMI 相關心因性休克
    \item Lactate $\geq$ 2.5 mmol/L
    \item SBP < 100 mmHg 或使用升壓劑
    \item LVEF < 45\%
    \item 休克發作至首次 Impella 支持 $\leq$ 24 小時
\end{enumerate}

\subsection{J-PVAD Registry 結果}

\begin{longtable}{p{5cm}p{4cm}p{4cm}}
\toprule
\textbf{族群} & \textbf{比例} & \textbf{30 天死亡率} \\
\midrule
\endfirsthead
\midrule
\textbf{族群} & \textbf{比例} & \textbf{30 天死亡率} \\
\midrule
\endhead
\bottomrule
\endfoot
符合 DanGer 標準的 STEMI-CS & 35.6\% & 37.6\% \\
\midrule
不符合標準的 STEMI-CS & - & 27.6\% \\
\midrule
院外心跳停止（OHCA） & - & 51.3\% \\
\midrule
機械性併發症（MC） & - & 39.8\% \\
\midrule
Non-STEMI-CS & - & 33.3\% \\
\bottomrule
\end{longtable}

\subsection{不符合標準患者的風險因子}

不符合 DanGer 標準但死亡率增加的因子：
\begin{itemize}
    \item 高齡
    \item 腎功能不全
    \item 使用 VA-ECMO
    \item 心室心律不整
    \item 敗血症
\end{itemize}

\note{真實世界中僅 1/3 的 AMI-CS Impella 患者符合 DanGer Shock 標準。對於不符合標準的亞群，需個人化風險評估來指導 Impella 使用。}

%============================================================
\section{長型 Troponin T 在 MI 診斷的應用}
%============================================================

\subsection{SuperTROPO 研究背景}

\begin{itemize}
    \item 現行高敏感度 cTnT 檢測測量完整與片段分子（total cTnT）
    \item 無法區分來源：急性 MI vs 慢性心肌損傷
    \item 新型檢測：僅測量完整或輕微片段的 cTnT（long cTnT）
\end{itemize}

\subsection{研究設計與結果}

\begin{longtable}{p{4cm}p{10cm}}
\toprule
\textbf{項目} & \textbf{內容} \\
\midrule
\endfirsthead
\midrule
\textbf{項目} & \textbf{內容} \\
\midrule
\endhead
\bottomrule
\endfoot
納入條件 & 急診 cTnT $\geq$ 14 ng/L 的連續患者 \\
\midrule
患者數 & 1,811 人（63.2\% 胸痛或呼吸困難）\\
\midrule
MI 患者 & 205 人（11.3\%），其中 148 人為 Type 1 MI \\
\midrule
關鍵發現 & Long cTnT 最低三分位（< 3.7 ng/L）僅 0.7\% 有 Type 1 MI \\
\bottomrule
\end{longtable}

\subsection{診斷效能}

\begin{center}
\begin{tabular}{lcc}
\toprule
& \textbf{Long cTnT} & \textbf{Total cTnT} \\
\midrule
AUC - Any MI & 0.833 & 0.782 \\
AUC - Type 1 MI & 0.839 & 0.777 \\
\bottomrule
\end{tabular}

\vspace{0.3cm}
兩者比較 P < .001
\end{center}

\subsection{臨床意義}

\begin{itemize}
    \item Long cTnT 在辨識 cTnT 升高患者中的 MI 表現優於 total cTnT
    \item 結合 long cTnT 與 total cTnT 可改善診斷準確度
    \item \textbf{但仍處於研究階段}，尚未準備納入常規臨床
    \item 未來可能結合 AI 與多重生物標記提升診斷
\end{itemize}

%============================================================
\section{本期重點總結}
%============================================================

\begin{enumerate}
    \item \textbf{PCI 後心絞痛}：需系統性評估，包括功能性造影與微血管評估

    \item \textbf{NSTE-ACS}：完整臨床與造影評估，多血管疾病應考慮完整血管重建

    \item \textbf{MI 定義}：通用定義有其價值，但臨床應用仍有挑戰

    \item \textbf{CHIP}：新興心血管風險因子，TET2 突變促進動脈粥狀硬化

    \item \textbf{心因性休克}：Impella 真實世界使用需謹慎，僅 1/3 符合 DanGer 標準

    \item \textbf{Troponin}：Long cTnT 有潛力改善 MI 診斷，但仍需更多研究
\end{enumerate}

%============================================================
\section{參考文獻}
%============================================================

\begin{enumerate}
    \item Crea F. Focus on acute coronary syndrome: mechanisms, diagnosis, management of cardiogenic shock, and new therapeutic targets. \href{https://doi.org/10.1093/eurheartj/ehag016}{\textit{Eur Heart J}. 2026;47:387-92.}

    \item Lombardi M, et al. Angina after percutaneous coronary interventions. \href{https://doi.org/10.1093/eurheartj/ehaf771}{\textit{Eur Heart J}. 2026;47:408-22.}

    \item Savonitto S, et al. Non-ST-segment elevation acute coronary syndrome: understanding the full spectrum to guide management. \href{https://doi.org/10.1093/eurheartj/ehaf958}{\textit{Eur Heart J}. 2026;47:423-37.}

    \item Mueller C, et al. Great debate: the universal definition of myocardial infarction is flawed and should be put to rest. \href{https://doi.org/10.1093/eurheartj/ehaf641}{\textit{Eur Heart J}. 2026;47:438-52.}

    \item von Scheidt M, et al. Clonal haematopoiesis of indeterminate potential and mortality in coronary artery disease. \href{https://doi.org/10.1093/eurheartj/ehaf602}{\textit{Eur Heart J}. 2026;47:453-69.}

    \item Arai R, et al. DanGer shock criteria and outcomes in acute myocardial infarction-related cardiogenic shock treated with Impella: the J-PVAD registry. \href{https://doi.org/10.1093/eurheartj/ehaf787}{\textit{Eur Heart J}. 2026;47:472-86.}

    \item Teppo K, et al. Long forms of cardiac troponin T for myocardial infarction diagnosis: the SuperTROPO study. \href{https://doi.org/10.1093/eurheartj/ehaf975}{\textit{Eur Heart J}. 2026;47:490-9.}

    \item Chen Z, et al. Arachidonic acid fuels inflammation by unlocking macrophage protein phosphatase 5 after myocardial infarction. \href{https://doi.org/10.1093/eurheartj/ehaf493}{\textit{Eur Heart J}. 2026;47:503-18.}
\end{enumerate}

\end{document}
